\documentclass{article}

\usepackage{Sweave}
\begin{document}
\Sconcordance{concordance:draftPaper.tex:draftPaper.Rnw:%
1 2 1 1 0 169 1}


\section{Introduction}
\subsection{Definition of masting}
\begin{enumerate}
\item Observations of masting in relation to seed predation and flowering
\end{enumerate}
\subsection{Major hypotheses for masting}
\begin{enumerate}
\item Predation satiation
\item Pollination coupling
\item Resource matching
\end{enumerate}
\subsection{Hypotheses operating at different reproductive stages like filters}
\begin{enumerate}
\item Pollination success (Pollination coupling)
\item Resources availability (Resource Matching)
\item Dispersal (Predation satiation)
\item Mechanisms are not mutually exclusive
\item Traits thus can be important to understand masting through their relative importance
\end{enumerate}
\subsection{Trait-based perspective}
\begin{enumerate}
\item Literature review on understanding masting from traits perspective 
\begin{itemize}
\item Larger seed is associate with higher CV
\item Animal dispersed seeds have higher CV than wind dispersed seeds
\item Wind pollination is related to higher CV
\item Recent papers also focus on resource perspective using leaf and shoot traits
\end{itemize}
\item Existing studies seem focus mainly on correlations
\item Not clear connection between traits and hypotheses and potential evolutionary drivers
\end{enumerate}
\subsection{Introduce framework linking hypothesis x traits x masting}
\begin{enumerate}
\item Link traits to hypotheses help understand drivers of masting for different species 
\item If the pattern is strong enough, we might be able to find it even using a binary masting definition
\item Here we use North American tree data to test this framework since they have information on their life-history
\begin{itemize}
\item Masting info available
\item Includes both masting and non-masting species
\item Not using CV data because don't have stand-level traits data
\end{itemize}
\end{enumerate}
\section{Methods}
\subsection{Data collection}
\begin{enumerate}
\item Data scraping
  \begin{itemize}
  \item Trait selection
  \item Data sources: Silvics of North America, Kew Seed Dataset, TRY Dataset
  \item Keywords used for data scraping
  \begin{itemize}
  \item Definition of strong masting vs. non-masting (referred to as masting vs. non-masting moving forward)
  \item Criteria for other traits
  \end{itemize}
  \end{itemize}
\item Data cleaning
  \begin{itemize}
  \item Cleaning of categorical traits
  \item Decision for numeric traits with a range
  \end{itemize}
\end {enumerate}
\subsection{Statistical analysis}
  \begin{enumerate}
  \item Main analysis for masting pattern and single trait: phyloglm
  \item Masting pattern with a multivariate analysis: NMDS
  \item Phylogenetic signal check for traits: Pagel's $\lambda$ and D-statistics
  \end{enumerate}
\section{Discussion}
\subsection{Summary of main findings}
\begin{enumerate}
\item Seed weight supports the predation satiation hypothesis
\item Pollination supports the pollination coupling hypothesis
\item These act at the start and the end of the reproductive cycle, provide evolutionary incentive
\item Discuss mechanisms in the context of existing literature
\begin{itemize}
\item Previous literature on seed mass and masting think from the resource depletion perspective
\item From evolutionary perspective, seeds attractive to animal dispersers will benefit more from masting
\end{itemize}
\item Other traits show some pattern but statistically non-significant
\item Transition to alternative hypotheses
\end{enumerate}
\subsection{Masting is a population-level phenomenon}
\begin{enumerate}
\item No evidence found at species level data does not mean traits are not related to masting
\item Examples of species that mast only under certain environmental conditions (E.g., \textit{Hymenaea coubaril})
\item Traits need to be measured locally
\end{enumerate}
\subsection{Different species selected under different hypotheses}
\begin{enumerate}
\item Introduce flowering masting vs. seed masting
\item Different species driven by different mechanisms
\item Need a more refined definition of masting and additional flowering-related traits
\item Transition to phylogeny
\end{enumerate}
\subsection{Evolutionary scale}
\begin{enumerate}
\item Angiosperm vs. gymnosperm results
\begin{itemize}
\item More masting gymnosperm sp
\item Longer leaf longevity in gymnosperm
\item More monoecious species in gymnosperm
\item Wind pollination is universal among gymnosperms
\item Gymnosperm and angiosperm have distinct trait space occupancy
\item Importance of looking for potential evolution support between gymnosperm and angiosperm
\begin{itemize}
\item Test if masting was gradually diluted in angiosperm as they transitioned to biotic pollination and faster life histories
\end{itemize}
\end{itemize}
\item Many traits have strong phylogenetic signal, it is hard to tease out traits from masting, which limits the inference
\item Data is strongly nested in phylogeny, and data is sparse.
\item Importance of testing within clades and collect more data (\textit{Quercus.})
\end{enumerate}
\end{document}
\subsection{Definetly include}
\begin{itemize}
\item Define masting
\item Introduce different hypotheses related to masting
\item Introduce that different hypotheses act on different stages of the reproductive process (e.g., flowering, pollination, seed maturation, dispersal)
\item Existing studies regarding functional traits and masting found: 
  \begin{itemize}
  \item High CV is common in species with high stem tissue density
  \item Short-lived leaves with low LMA is associate with low CV
  \item Large seed is associate with high CV
  \item floral transition is associate with year-to-year variability
  \item wind pollination is associate with hight CV
  \end{itemize}
\item Existing studies lack emphasis on the role of reproductive traits as mechanistic links between masting and hypotheses.
\item Especially on relationship between seed mass and masting focus on an ecological constraint perspective, where seed mass reflects the cost of reproduction. However, we try to capture the evolutionary mechanisms underlying masting benefits. For masting to be adaptive under the predator satiation hypothesis, seeds must be consumed by animals. Therefore, fruit size might be more relevant than seed mass alone.
\item Our study propose a framework connecting:
\begin{itemize}
\item reproductive traits
\item stages of reproduction
\item specific masting hypotheses
\end{itemize}

\item This framework can clarify which traits are under selection in masting species


\end{itemize}

\section{Discussion}
\subsection{Definetly include}
\begin{itemize}
\item Seed mass and wind pollination are reproductive traits significantly associated with masting patterns. This might suggest that reproductive allocation plays a key role in shaping masting behavior, supporting predation satiation hypothesis and pollination coupling hypothesis. 
\item However, the fact that we didn't find relationships with other traits indicates that trait-based explanations alone are not sufficient to fully explain masting.
\item \textbf{Gymnosperm are more likely to be masting species}, (is there a possibility that masting might be an ancestral strategy that becomes lost in some angiosperm lineages later?)
\item We found moderate phylogenetic signal for masting in both gymnosperms and angiosperms, which align with other people's finding (e.g.,: Pearse et al., 2020, found a modest phylogenetic signal in CVp), then link to next point.
\item \textbf{Masting is strongly related to environment as well}, some species might be exhibit masting or not depending on the environment they are in, broad trait-based analysis do not capture environmental drivers, and that's the reason why we can't use any masting metrics capturing other characteristics of masting.
\item Different hypotheses may act at different stages (for flower masting and seed masting species), playing the selection on different species, need better data on flowering traits and finer resolution on masting definition.
\item Existence of clear counterexamples especially within the clade. Focusing on within-clade variation (oaks) may be more informative. Need more detailed monitoring on masting pattern and reproductive traits for more species in oak family, potentially with data for same species in different environment.
\end{itemize}
